%% kapa: exam=1 round=34 year=2023 period=1 form=A subject=economics subject_ko=경제학원론 seq=16 number=56
%% source: Q-Net 감정평가사 기출문제 (cst00309) artlSeq=5211789 file=raw/1008/5211789/2023년도 제34회 감정평가사 1차시험 1교시.pdf
%% conversion: kordoc 4.24.51 -> tools/gen1.py (auto); review=proofread-sonnet; flags=none
\begin{question}{56}
외부성에 관한 설명으로 옳은 것은?
\begin{choices}
  \item 생산의 부정적 외부성이 있는 경우 사회적 최적생산량이 시장균형생산량보다 크다.
  \item 생산의 부정적 외부성이 있는 경우 사적 한계비용이 사회적 한계비용보다 작다.
  \item 소비의 부정적 외부성이 있는 경우 사적 한계편익이 사회적 한계편익보다 작다.
  \item 코즈(R. Coase)의 정리는 거래비용의 크기와 무관하게 민간경제주체들이 외부성을 스스로 해결할 수 있다는 정리를 말한다.
  \item 공유자원의 비극(tragedy of the commons)은 긍정적 외부성에서 발생한다.
\end{choices}
\end{question}
\answer{2}
